\documentclass[hyperref={colorlinks=true,linkcolor=blue}, aspectratio=169]{beamer}

\usepackage{animate}

\usepackage{booktabs}
\setbeamertemplate{caption}[numbered]
\setbeamertemplate{footline}[frame number]
\setbeamertemplate{navigation symbols}{}

\usepackage{amsmath}
\newcommand*{\norm}[1]{\left\lVert#1\right\rVert}
\newcommand{\argmin}{\arg\!\min}

\usepackage{color,soul}
\usepackage{colortbl}

\setbeamerfont{footnote}{size=\tiny}

\usepackage{listings,lstautogobble}
\lstset{language=Python,
	basicstyle=\scriptsize\ttfamily,
	commentstyle=\ttfamily\itshape\color{gray},
	showstringspaces=false,
	autogobble=true
}

\usepackage{caption}
\captionsetup{justification=centering,margin=1cm}
\usepackage{subcaption}

\usepackage{multirow}

\usepackage{siunitx}
\usepackage{tikz}
\tikzstyle{process} = [rectangle, minimum width=3cm, minimum height=1cm, text centered, draw=black, fill=orange!30]
\usetikzlibrary{shapes,positioning}
\usepackage{ulem}

\AtBeginSection[]{
	\begin{frame}
		\vfill
		\centering
		\begin{beamercolorbox}[sep=8pt,center,shadow=true,rounded=true]{title}
			\usebeamerfont{title}\insertsectionhead\par%
		\end{beamercolorbox}
		\vfill
	\end{frame}
}

\addtobeamertemplate{footnote}{\hskip -2em}{}

\usepackage[capitalise,noabbrev]{cleveref}
\usepackage[numbers,sort&compress]{natbib}
\bibliographystyle{unsrt}

\setbeamercolor{block title}{fg=black,bg=yellow}

\title{In vivo 7~T diffusion MRI: \\High spatio-angular-temporal resolution pursuit}
\author[Z Tan]{Zhengguo Tan \inst{1}, Patrick A Liebig \inst{2}, Robin M Heidemann \inst{2}, Frederik B Laun \inst{3}, Florian Knoll \inst{1}}
\institute{\inst{1} Artificial Intelligence in Biomedical Engineering, University of Erlangen-Nuremberg, Erlangen, Germany \\
\inst{2} Siemens AG, Erlangen, Germany \\
\inst{3} Institute of Radiology, University Hospital Erlangen, Erlangen, Germany}
\date{07 June 2023}

\begin{document}
	
	\frame{\titlepage}

	% ================= %	
	\section{Motivation: In vivo Diffusion MRI at 7~T}
	
	\begin{frame}{MUSE \footnote{Chen NK, et al. A robust multi-shot scan strategy for high-resolution diffusion weighted MRI enabled by multiplexed sensitivity-encoding (MUSE). \textit{NeuroImage} (2013). doi:~\href{https://doi.org/10.1016/j.neuroimage.2013.01.038}{10.1016/j.neuroimage.2013.01.038}} faces challenges in ultra-high field}
		\begin{itemize}
			\item [\checkmark] MUSE: self-gated multi-shot interleaved EPI (iEPI) DWI reconstruction
		\end{itemize} 
		\begin{figure}
			\includegraphics[width=0.7\textwidth]{figures/4-shot.png}
		\end{figure}
		\vspace{0.5em}
		\begin{itemize}
			\item [$\star$] <1-> requires \uline{fully-sampled in-plane acquisition}
			\item [$\star$] <2-> only limited number of shots have been reported
			\item [$\star$] <3-> increases specific absorption rate (SAR) at ultra-high field
		\end{itemize}
	\end{frame}

	\begin{frame}{Undersampled iEPI Acquisition and Joint $k$-$q$-Space Reconstruction}
	This work presents preliminary results on 
	\begin{enumerate}
		\item undersampled iEPI acquisition to reduce scan time and to reduce SAR
		\item joint $k$-$q$-space reconstruction to boost SNR
	\end{enumerate}
	\end{frame}

	% ================= %
	\section{Methods}
	
	\begin{frame}{Undersampled iEPI with Shift Encoding}
		\begin{figure}
			\includegraphics[width=0.7\textwidth]{figures/2-shot.png}
		\end{figure}
		\begin{itemize}
			\item [(A)] <1-> Given the acceleration along $k_y$ as $R$ $\rightarrow$ $(R \times N_\mathrm{shot})$-fold undersampling per shot
			\item [(B)] <2-> $k_y$ shift encoding along diffusion encoding $\rightarrow$ complementary $k$-$q$-space sampling
		\end{itemize}
	\end{frame}

	\begin{frame}{In Vivo Acquisition}
		\begin{columns}
			\begin{column}{0.6\textwidth}
				\begin{itemize}
					\item three-shell diffusion sampling
					\item \SI{1.4}{mm} isotropic resolution and 68 slices
					\item $k_y$ undersampling $R = 3$, partial Fourier $6/8$
					\vspace{1em}
					\item 1- and 2-shot scans
				\end{itemize}
			\end{column}
			\begin{column}{0.4\textwidth}
				\begin{figure}
					\includegraphics[width=\columnwidth]{figures/spheres_b2500.png}
				\end{figure}
			\end{column}
		\end{columns}
	\end{frame}

	\begin{frame}{Joint $k$-$q$-Space Reconstruction}
		Shot phase variation informed joint reconstruction:
		
		\begin{equation}
			\mathrm{argmin}_{\mathbf{x}_q} \left\lVert \mathbf{y}_{c,q,s} - \mathbf{P}_{q,s} \mathbf{F} \mathbf{S}_c \{ e^{i \mathbf{\Phi}_{q,s}} \mathbf{x}_{q} \} \right\rVert_2^2 + \lambda \mathbf{R} (\mathbf{x}_q)
		\end{equation}
	
		\begin{equation}
			\begin{aligned}
				\mathbf{x} &: \text{diffusion weight images}, \\
				\mathbf{\Phi} &: \text{shot-to-shot phase variation}, \\
				\mathbf{S} &: \text{coil sensitivity maps}, \\
				\mathbf{F} &: \text{2D Fourier transform}, \\
				\mathbf{P} &: \text{sampling mask}, \\
				\mathbf{y} &: \text{measured raw data}.
			\end{aligned}
		\end{equation}
		\vspace{0.5em}
		Indices: $q$ - diffusion encoding; $s$ - shot; $c$ - coil
	\end{frame}

	\begin{frame}{Overlapping Locally Low Rank (LLR) Regularization}
		\begin{columns}
			\begin{column}{0.4\textwidth}
				\begin{enumerate}
					\item Stride set as 1, leading to overlapping local patches
					\item Use local patches to create spatial-contrast (diffusion) matrices
					\item singular value thresholding (SVT) \footnote{Cai JF, Candes EJ, Shen Z. A singular value thresholding algorithm for matrix completion. \textit{SIAM J Optim} (2010). doi: \href{https://doi.org/10.1137/080738970}{10.1137/080738970}} on these matrices
					\item Backpass (adjoint) the SVT matrices back to full image matrices
				\end{enumerate}
			\end{column}
			\begin{column}{0.6\textwidth}
				\animategraphics[autoplay,loop,controls,width=\columnwidth]{1}{figures/llr/frame_}{0}{5}
			\end{column}
		\end{columns}
	\end{frame}

	% ================= %
	\section{Results}
	
	\begin{frame}{1-Shot EPI: LLR Regularization Strength}
		\begin{figure}
			\includegraphics[width=0.8\textwidth]{figures/results/Pat3_PI_LLR_regu.png}
		\end{figure}
	\end{frame}

	\begin{frame}{1-Shot EPI: Overlapping LLR}
		\begin{figure}
			\includegraphics[width=0.9\textwidth]{figures/results/DWI_LLR_OverlapBlocks.png}
		\end{figure}
	\end{frame}

	\begin{frame}{2-Shot EPI: Parallel Imaging Reconstruction}
		\begin{figure}
			\includegraphics[width=0.75\textwidth]{figures/results/fig3_a.png}
		\end{figure}
	\end{frame}

	\begin{frame}{2-Shot EPI: Overlapping LLR Regularization}
		\begin{figure}
			\includegraphics[width=0.75\textwidth]{figures/results/fig3_b.png}
		\end{figure}
	\end{frame}


	% ================= %
	\section{Conclusion}
	
	\begin{frame}{Conclusion}
		
		\begin{itemize}
			\item <1-> Undersampled iEPI acquisition at ultra-high field
			\vspace{1em}
			\item <2-> Joint $k$-$q$-space reconstruction
			\item <2-> Spatial-diffusion overlapping LLR regularization
			\vspace{1em}
			\item <3-> Preliminary in vivo results show better image quality compared to MUSE
			\vspace{1em}
			\item <4-> Please refer to: \\ \url{https://github.com/ZhengguoTan/demo_jets_diffusion_mri_7t.git} \\
			for future updates!
			\item <5-> MR-Pub III at Reproducible Research Study Group at 13:45 on Thursday
			\vfill
			\item <6-> \colorbox{yellow}{Thank You for Your Attention!}
		\end{itemize}
	\end{frame}

\end{document}